\documentclass[11pt,a4paper]{scrartcl}

\usepackage[lithuanian]{babel}
\usepackage{../../1.\space Vienanariai\space ir\space daugianariai/manostilius_lt}
\usepackage{tasks}
\renewcommand{\theproblem}{\arabic{problem}}

\title{Trečio ir ketvirto laipsnio lygčių sprendimas}
\author{Andrius Berniukevičius}

\begin{document}

\maketitle
\section*{Kaip dirbti su uždaviniais?}

Pirmiausia kiekvieną uždavinį pabandykite išspręsti savarankiškai. Neskubėkite ieškoti užuominų — skirkite laiko apgalvoti, ką jau žinote ir kokį sprendimo būdą galėtumėte pritaikyti. Jei įstrigote, atsiverskite atitinkamą „Hintai 1“ užuominą. Pagalvojate. Jei jos nepakanka, tik tada skaitykite „Hintai 2“, o prireikus — „Hintai 3“. Užuominas naudokite palaipsniui, nežiūrėkite jų iš anksto.

Paties pagimdytas sprendimas yra šimtą kartų vertingesnis už įvaikintą sprendimą.

\section{Uždaviniai}

\begin{problem}
Išspręskite kubines lygtis realiųjų skaičių aibėje:
\begin{tasks}[label={(\arabic*)}, label-offset=0.9em](2)
    \task $x^3 + 5x^2 + 6x = 0$
    \task $5y^3 - 8y^2 + 3y = 0$
    \task $-3a^3 + a^2 +2a = 0$
    \task $-m^3 + 4m^2 +3m = 0$
\end{tasks}
\end{problem}

\begin{problem}
    Raskite daugianario $P(x) = x ^{3} +5x ^{2} -4x-2$ realiąsias šaknis, jei žinoma, kad daugianarį galima užrašyti pavidalu:
     $$P(x) = (x-1)(x ^{2} +bx+c) $$    
\end{problem}

\begin{problem}
    Raskite daugianario $P(x) =x^3+4 x^2+x-6$ realiąsias šaknis, jei žinoma, kad daugianarį galima užrašyti pavidalu:
     $$P(x) = (x+3)(x ^{2} +bx+c) $$    
\end{problem}

\begin{problem}
Išspręskite ketvirtojo laipsnio lygtis realiųjų skaičių aibėje:
\begin{tasks}[label={(\arabic*)}, label-offset=0.9em](1)
    \task $(2x ^{2} +3x-5) ^{2} -(x ^{2} +4x-3) ^{2} = 0$
    \task $(3x+2) ^{2} = (2x ^{2} +2x+1) ^{2} $
    \task $(x ^{2} -3x+1) ^{2} +2(x-3)(x ^{2} -3x+1)+x ^{2} -6x+9 = 0$
    \task $(x ^{2} -4) ^{2} -(2x ^{2} -8)(2x ^{2} +x)+4x ^{4}+4x ^{3} +x ^{2} = 0 $
    \task $6x ^{2} (x ^{2} +3x+2)-3x(x ^{2} +3x+2)+ 2(x ^{2} +3x+2) = 0$
    \task $(2x+5)(x^2 + 5x + 6)+(3x+10)(x ^{2} +5x+6) = 0$ 
    \task $(x+1)(x ^{2} +1)(x-1) = 15$
\end{tasks}
\end{problem}

\newpage
\section{Atsakymai}

\begin{enumerate}[label=\textbf{\arabic*.}, leftmargin=*]
    \item
    \begin{enumerate}[label={(\arabic*)}]
        \item $x\in\{-3,-2,0\}$
        \item $y\in\left\{0,\frac35,1\right\}$
        \item $a\in\left\{-\frac23,0,1\right\}$
        \item $m\in\{0,2-\sqrt7,2+\sqrt7\}$
    \end{enumerate}
    \item
    $x\in\{1,-3+\sqrt7,-3-\sqrt7\}$
    \item
    $x\in\{-3,-2,1\}$
    \item
    \begin{enumerate}[label={(\arabic*)}]
        \item $x\in\left\{-1,2,\frac{-7-\sqrt{145}}6,\frac{-7+\sqrt{145}}6\right\}$
        \item $x\in\left\{-\frac32,-1,-\frac12,1\right\}$
        \item $x\in\{1-\sqrt3,1+\sqrt3\}$
        \item Sprendinių realiųjų skaičių aibėje nėra.
        \item $x\in\{-2,-1\}$
        \item $x\in\{-3,-2\}$
        \item $x\in\{-2,2\}$
    \end{enumerate}
\end{enumerate}

\newpage
\section{Hintai 1}

\begin{enumerate}[label=\textbf{\arabic*.}, leftmargin=*]
    \item Iškelkite bendrąjį daugiklį už skliaustų.
    \item Atskliauskite duotą sandaugą ir palyginkite gauto daugianario koeficientus su duotaisiais.
    \item Atskliauskite duotą sandaugą ir palyginkite gauto daugianario koeficientus su duotaisiais.
    \item Pirmose keturiose lygtyse pamėginkite pritaikyti kvadratų skirtumo formulę, iškelti bendrąjį daugiklį arba atpažinti pilnąjį kvadratą. Likusiose iškelkite pasikartojančius daugiklius.
    \begin{enumerate}[label={(\arabic*)}]
        \item Pritaikykite $A^2-B^2=(A-B)(A+B)$.
        \item Perkelkite viską į vieną lygties pusę ir vėl pritaikykite kvadratų skirtumo formulę.
        \item Patikrinkite, ar kairioji pusė yra dviejų vienodų reiškinių suma, pakelta kvadratu.
        \item Pabandykite užrašyti reiškinį kaip $A^2-2AB+B^2$.
        \item Iškelkite bendrą daugianarį $x^2+3x+2$.
        \item Iškelkite bendrą daugianarį $x^2+5x+6$.
        \item Kairėje pusėje sujunkite daugiklius $x+1$ ir $x-1$.
    \end{enumerate}
\end{enumerate}

\newpage
\section{Hintai 2}

\begin{enumerate}[label=\textbf{\arabic*.}, leftmargin=*]
    \item Iškelkite kintamąjį prieš skliaustus, kad liktų kvadratinis trinaris. Sandauga lygi nuliui, kai bent vienas iš jos dauginamųjų lygus nuliui.
    \begin{enumerate}[label={(\arabic*)}]
        \item $x(x^2+5x+6)=0$.
        \item $y(5y^2-8y+3)=0$.
        \item $a(-3a^2+a+2)=0$.
        \item $m(-m^2+4m+3)=0$.
    \end{enumerate}
        \item Atskliaudus gauname $x^3+(b-1)x^2+(c-b)x-c$.
        \item Atskliaudus gauname $x^3+(b+3)x^2+(c+3b)x+3c$.

    \item
    \begin{enumerate}[label={(\arabic*)}]
        \item Pažymėję $A=2x^2+3x-5$, $B=x^2+4x-3$, gausite $(A-B)(A+B)=0$.
        \item Pažymėję $A=3x+2$, $B=2x^2+2x+1$, gausite $(A-B)(A+B)=0$.
        \item Pažymėję $A=x^2-3x+1$, $B=x-3$, gausite $(A+B)^2=0$.
        \item Pažymėję $A=x^2-4$, $B=2x^2+x$, gausite $(A-B)^2=0$.
        \item Pažymėję $A=x^2+3x+2$, gausite $A(6x^2-3x+2)=0$.
        \item Pažymėję $A=x^2+5x+6$, gausite $A(5x+15)=0$.
        \item Pažymėję $A=x^2-1$, gausite $A(A+2)=15$.
    \end{enumerate}
\end{enumerate}

\newpage
\section{Hintai 3}

\begin{enumerate}[label=\textbf{\arabic*.}, leftmargin=*]
    \item Išspręskite gautas tiesines arba kvadratines lygtis.
    \begin{enumerate}[label={(\arabic*)}]
        \item Sprendiniai: $x=-3,-2,0$.
        \item Sprendiniai: $y=0,\frac35,1$.
        \item Sprendiniai: $a=-\frac23,0,1$.
        \item Sprendiniai: $m=0$ arba $m=2\pm\sqrt7$.
        \end{enumerate}
        \item $b=6$, $c=2$, todėl $(x-1)(x^2+6x+2)=0$; šaknys: $1,-3\pm\sqrt7$.
        \item $b=1$, $c=-2$, todėl $(x+3)(x^2+x-2)=0$; šaknys: $-3,-2,1$.

    \item Įstatykite iš „Hintai 2“ gautas išraiškas ir išspręskite lygtis.
    \begin{enumerate}[label={(\arabic*)}]
        \item Pažymėję $A-B=x^2-x-2$, $A+B=3x^2+7x-8$, gausite: $$(x-2)(x+1)(3x^2+7x-8)=0.$$ 
        Sprendiniai $x=2,-1,\frac{-7\pm\sqrt{145}}6$.
        \item Pažymėję $A-B=-2x^2+x+1$, $A+B=2x^2+5x+3$, gausite:
        $$(2x+1)(x-1)(2x+3)(x+1)=0.$$
        Sprendiniai $x=-\frac12,1,-\frac32,-1$.

        \item Pažymėję $A+B=x^2-2x-2$, gausite:
        $$(x^2-2x-2)^2=0.$$
        Sprendiniai $x=1\pm\sqrt3$.

        \item Pažymėję $A-B=-x^2-x-4$, gausite:
        $$(x^2+x+4)^2=0.$$
        Kadangi šio trinario diskriminantas neigiamas, realiųjų sprendinių nėra.

        \item Pažymėję $A=x^2+3x+2$, gausite:
        $$A(6x^2-3x+2)=0.$$
        Antrasis daugiklis neturi realiųjų šaknų, todėl sprendiniai $x=-1,-2$.

        \item Pažymėję $A=x^2+5x+6$, gausite:
        $$A(5x+15)=0.$$
        Išskaidę daugiklius, gausite $(x+2)(x+3)^2=0$. Sprendiniai $x=-2,-3$.

        \item Pažymėję $A=x^2-1$, gausite:
        $$A(A+2)=15.$$
        Kadangi $A(A+2)=x^4-1$, gausite $x^4=16$. Realieji sprendiniai $x=-2,2$.
    \end{enumerate}
\end{enumerate}

\end{document}
